\documentclass[12pt]{article}
\usepackage[utf8]{inputenc}
\usepackage[T1]{fontenc}
\usepackage{hyperref}
\title{Mitla Zapotec Language Resources for K-12 Curriculum Development}
\author{Compiled by Claude Sonnet 5.5 (claude-sonnet-5-5) for Andr\'as Kornai}
\date{}
\begin{document}
\maketitle

\section{Introduction}
Mitla Zapotec (ISO 639-3 zaw, Glottocode mitl1236) is an Otomanguean language of the Oaxaca valley, Mexico, with about 20,000 speakers. Only a few reference works turned up, mainly from SIL, and nothing aimed at schools was found.

\section{Primary Educational Resources}
No primers, textbooks or school programmes were found in the search results.

\section{Linguistic and Reference Materials}
The main dictionary is Carol and Morris Stubblefield's \emph{Diccionario Zapoteco de Mitla, Oaxaca}, published by the Instituto Ling\"u\'istico de Verano. The first edition is from 1991 and a second from 2021 (xiii+146 pages), with Spanish as the metalanguage. The same authors compiled \emph{Mitla Zapotec texts} (1994). Briggs's 1961 \emph{Mitla Zapotec grammar} is the standard grammatical reference. No online copy was found. The links below are bibliographic or database entries.
\begin{itemize}
\item Glottolog reference, Stubblefield and Stubblefield 1991: \url{https://glottolog.org/resource/reference/id/160564}
\item Glottolog reference, Stubblefield and Stubblefield 2021: \url{https://glottolog.org/resource/reference/id/668076}
\item WALS reference record, Stubblefield and Stubblefield 1991: \url{https://wals.info/refdb/record/Stubblefield-and-Stubblefield-1991}
\item Oto-Manguean database entry: \url{https://oto-manguean.surrey.ac.uk/Info/ZAW}
\item Glottolog languoid page: \url{https://glottolog.org/resource/languoid/iso/zaw}
\item Wikipedia overview: \url{https://en.wikipedia.org/wiki/Mitla_Zapotec}
\end{itemize}

\section{Religious and Community Resources}
None specific to Mitla Zapotec were found in the search results.

\section{Digital and Audio Resources}
No digital dictionary, audio or video resources were found in the search results.

\section{Curriculum Development Context}
The language is reported as stable, with all members of the community using it as a first language, but it is not known to be taught in schools. The Stubblefield dictionary, texts and Briggs grammar would be the first sources for a curriculum. Their availability and licensing would have to be checked, and work with the community would be needed to produce school materials.

\end{document}